% answers_short.tex —— 2015 年·湖北卷·文科 · 简短答案速查
% 由 tools/make_short_answers.py 自 answers.tex 的 \gksolhead/\gkanswer 提取，
% 并与 2015/2015_数学真题_湖北文_答案要点.md 逐题比对后生成，勿手改（改源 MD 后重跑）。

\begin{gknotice}
\textbf{使用说明}：本节只给最终答案便于快速核对，完整推导见解析版。答案为编写者独立推导，\textbf{非官方标准答案}。
\end{gknotice}

\gkpart{选择题}
\begin{center}\begin{tabular}{ccccccccccc}
\toprule
题号 & 1 & 2 & 3 & 4 & 5 & 6 & 7 & 8 & 9 & 10 \\
\midrule
答案 & B & B & C & C & A & C & D & D & D & C \\
\bottomrule
\end{tabular}\end{center}

\gkpart{填空题}
\begin{itemize}[leftmargin=2.6em]
  \item[\textbf{11.}] \(9\)
  \item[\textbf{12.}] \(10\)
  \item[\textbf{13.}] \(2\)
  \item[\textbf{14.}] (1) \(a=3\)；(2) \(6000\) 人
  \item[\textbf{15.}] \(100\sqrt6\)
  \item[\textbf{16.}] (1) \((x-1)^2+(y-\sqrt2)^2=2\)；(2) \(-(1+\sqrt2)\)
  \item[\textbf{17.}] \(2\sqrt2-2\)（此时 \(g(a)\) 的最小值为 \(3-2\sqrt2\)）。
\end{itemize}

\gkpart{解答题}
\begin{itemize}[leftmargin=2.6em]
  \item[\textbf{18.}] (1) 表格补全为 \(x=\dfrac{\pi}{12},\dfrac{\pi}{3},\dfrac{7\pi}{12},\dfrac{5\pi}{6},\dfrac{13\pi}{12}\)（对应值 \(0,5,0,-5,0\)）；\(f(x)=5\sin\left(2x-\dfrac{\pi}{6}\right)\)；(2) \(\left(-\dfrac{\pi}{12},0\right)\)
  \item[\textbf{19.}] (1) 当 \(d=2\) 时 \(a_n=2n-1\)，\(b_n=2^{n-1}\)；当 \(d=\dfrac29\) 时 \(a_n=\dfrac{2n+79}{9}\)，\(b_n=9\left(\dfrac29\right)^{n-1}\)；(2) \(T_n=6-\dfrac{2n+3}{2^{n-1}}\)
  \item[\textbf{20.}] (1) 成立（\(DE\perp\) 平面 \(PBC\)）；\(EBCD\) 是鳖臑，四个面的直角分别为 \(\angle BCD\)、\(\angle BCE\)、\(\angle CED\)、\(\angle BED\)；(2) \(4\)
  \item[\textbf{21.}] (1) \(f(x)=\dfrac{\mathrm{e}^x-\mathrm{e}^{-x}}{2}\)，\(g(x)=\dfrac{\mathrm{e}^x+\mathrm{e}^{-x}}{2}\)；当 \(x>0\) 时 \(f(x)>0\)，\(g(x)>1\) 成立；(2) 成立（由 \(\dfrac{f(x)}{x}\) 介于 \(1\) 与 \(g(x)\) 之间，结合 \(a\leqslant0\)、\(b\geqslant1\) 放缩即得）
  \item[\textbf{22.}] (1) 曲线 \(C\) 的方程为 \(\dfrac{x^2}{16}+\dfrac{y^2}{4}=1\)；(2) 存在最小值，最小值为 \(8\)
\end{itemize}
